\documentclass[11pt,a4paper]{article}

\usepackage[utf8]{inputenc}
\usepackage[T1]{fontenc}
\usepackage[english]{babel}
\usepackage[a4paper,margin=2.4cm]{geometry}

\usepackage{amsmath, amssymb, amsthm}
\usepackage{mathtools}
\usepackage{booktabs}
\usepackage{tabularx}
\usepackage{array}
\usepackage{caption}
\usepackage{float}
\usepackage{enumitem}
\usepackage{microtype}
\usepackage{xcolor}
\usepackage{listings}
\usepackage{fancyhdr}
\usepackage{titlesec}
\usepackage{tikz}
\usetikzlibrary{positioning, arrows.meta, fit, backgrounds, calc, shapes.geometric}
\usepackage[hidelinks]{hyperref}

%  ---  Palette shared with the companion report  ---
\definecolor{ink}{RGB}{26,30,36}
\definecolor{slate}{RGB}{92,101,112}
\definecolor{mist}{RGB}{233,235,238}
\definecolor{rulegrey}{RGB}{188,194,201}
\definecolor{steel}{RGB}{47,72,102}
\definecolor{steellight}{RGB}{124,146,172}
\definecolor{signal}{RGB}{160,32,45}

\hypersetup{
    colorlinks=true,
    linkcolor=steel,
    citecolor=steel,
    urlcolor=steel,
    pdftitle={The Six-Room Survey: An EPDDL Domain Executed on ePlanSys},
    pdfauthor={Haniel Ulises},
}

\color{ink}

\setlength{\headheight}{22pt}
\pagestyle{fancy}
\fancyhf{}
\fancyhead[L]{\small\textcolor{slate}{\itshape The Six-Room Survey}}
\fancyhead[R]{\small\textcolor{slate}{\thepage}}
\renewcommand{\headrulewidth}{0.4pt}
\renewcommand{\headrule}{\hbox to\headwidth{\color{rulegrey}\leaders\hrule height \headrulewidth\hfill}}

\titleformat{\section}{\normalfont\large\bfseries\color{ink}}{\thesection}{0.7em}{}
\titleformat{\subsection}{\normalfont\normalsize\bfseries\color{ink}}{\thesubsection}{0.7em}{}
\titleformat{\paragraph}[runin]{\normalfont\bfseries\color{slate}}{}{0em}{}[.]

\captionsetup{font=small, labelfont={bf,color=slate}, textfont=footnotesize}

\theoremstyle{definition}
\newtheorem{definition}{Definition}

\lstdefinestyle{sys}{
    backgroundcolor=\color{mist},
    basicstyle=\ttfamily\scriptsize\color{ink},
    keywordstyle=\color{steel}\bfseries,
    commentstyle=\color{slate}\itshape,
    stringstyle=\color{slate},
    morecomment=[l]{\#},
    morekeywords={executor,ros__parameters,bt_builder_plugin,bt_node_plugins,
        planner,plan_solver_plugins,EPISTEMIC,conditional_plan,action_mapping,
        Sequence,ReactiveSequence,CheckKnowledge,ApplyEpistemicUpdate,
        EpistemicSwitch,CheckEpistemicGoal,ExecuteAction,true,false,ok,refused},
    frame=single,
    rulecolor=\color{rulegrey},
    framesep=5pt,
    xleftmargin=6pt,
    xrightmargin=0pt,
    breaklines=true,
    showstringspaces=false,
    columns=fullflexible,
    captionpos=b,
}
\lstset{style=sys}

%  ---  Notation  ---
\newcommand{\Ag}{\mathcal{A}}
\newcommand{\Mod}{\mathcal{M}}
\newcommand{\Ev}{\mathcal{E}}
\newcommand{\prd}{\otimes}
\newcommand{\Kop}[1]{K_{#1}}
\newcommand{\Kw}[1]{\mathit{Kw}_{#1}}
\newcommand{\code}[1]{\texttt{\small #1}}

\begin{document}

\begin{titlepage}
\centering
\vspace*{2.4cm}
{\color{rulegrey}\rule{\textwidth}{0.8pt}}\\[0.9cm]
{\LARGE\bfseries The Six-Room Survey\par}
\vspace{0.5cm}
{\large\color{slate} An EPDDL Domain Where the Map Supplies the Knowledge,\\ Executed on ePlanSys\par}
\vspace{0.7cm}
{\color{rulegrey}\rule{\textwidth}{0.8pt}}\\[1.4cm]

{\large Haniel Ulises\par}
\vspace{0.3cm}
{\color{slate} Epistemic Robotics}\\
\vspace{0.15cm}
{\color{slate}\today}

\vfill
{\color{slate}\small
A single robot, six rooms, and a goal that no classical planner can state:
not to reach the east wing, but to \emph{know} whether it is passable.
The observation is not imagined by the planner --- it is produced by
\code{plansys2\_epistemic\_perception} from an occupancy grid the robot's own
laser builds while it drives.}
\vspace{1.2cm}
\end{titlepage}

\section{What the domain has to express}

A survey mission over a building is a poor fit for classical planning, and the
reason is not the geometry. A robot can be told to drive to the far side of a
building and a metric planner will route it there. What it cannot be told is
that the mission succeeds only once the robot \emph{knows} whether the far side
is passable --- a condition about the robot's information, not about the world.

The distinction has teeth when the answer is not determined in advance. Rooms
nobody has entered are neither known to be clear nor known to be blocked, and a
plan that assumes either is unexecutable in half of the situations it will meet.
What the mission needs is a policy: drive somewhere the question can be
answered, answer it, and continue accordingly.

This report describes such a domain --- \code{building-rooms} --- its solution,
and its execution on ePlanSys, together with the defect that currently stops the
execution half short.

\section{The domain}

The building has six rooms. Rooms~3 and~6 lie east of a cross corridor and are
the subject of the survey; the robot starts in room~1 at the west end. Four
propositions suffice:

\begin{center}
\begin{tabular}{@{}ll@{}}
\toprule
\code{blocked-r3}, \code{blocked-r6} & whether each east room is obstructed \\
\code{at-door3 ?i}, \code{at-door6 ?i} & whether \code{?i} stands where that room can be seen into \\
\bottomrule
\end{tabular}
\end{center}

Two ontic actions drive to the doorways, and are public: in a fleet every robot
tracks the others' positions, so the effect is common knowledge. Two sensing
actions look in. Their type is what carries the epistemic content:

\begin{lstlisting}
(:action inspect3
  :parameters (?i - agent)
  :action-type (semi-private-sensing
    (e-inspect3-blocked ?i)
    (e-inspect3-clear ?i) )
  :observability-conditions
    (:and
      (?i Fully)
      (default Partially) ) )
\end{lstlisting}

Semi-private sensing says that the agent which looks learns the answer, while
anyone else present sees it look without learning what it saw. With one robot
the distinction is invisible; modelling it as fully public would nevertheless be
a lie, and one that stops being harmless the moment a second robot is added.

Each sensing action carries two events, one per answer, and neither is known to
apply when the plan is built. That is precisely what makes the solution branch.

\subsection{The instance}

What the initial state does \emph{not} say is whether either room is
obstructed. The theory leaves both open, so the model designates four worlds
rather than one --- a building nobody has surveyed is genuinely undetermined:

\begin{lstlisting}
(:init
  (:and
    (:forall (?i - agent) ([C. All] (not (at-door3 ?i))))
    (:forall (?i - agent) ([C. All] (not (at-door6 ?i))))
    (:forall (?i - agent) ([C. All] (<Kw. ?i> (blocked-r3))))
    (:forall (?i - agent) ([C. All] (<Kw. ?i> (blocked-r6)))) ) )

(:goal (and ([Kw. r1] (blocked-r3)) ([Kw. r1] (blocked-r6))))
\end{lstlisting}

The goal is $\Kw{r_1}\,\mathit{blocked}\text{-}r_3 \wedge
\Kw{r_1}\,\mathit{blocked}\text{-}r_6$: knowing \emph{whether}, not knowing
\emph{that}. A mission that succeeds only when the rooms turn out to be clear is
not a mission, and a plan for it would be unexecutable in the half of the worlds
where they are not.

\section{Grounding and search}

Grounding with \textsc{plank} yields one agent, four atoms, four grounded
actions and an initial model of four worlds, all designated:

\begin{lstlisting}
agents: ['r1']
atoms:  ['at-door3_r1', 'blocked-r6', 'at-door6_r1', 'blocked-r3']
actions: ['goto-door3_r1', 'goto-door6_r1', 'inspect3_r1', 'inspect6_r1']
initial worlds: 4   designated: 4
\end{lstlisting}

\textsc{Aletheia} selects its heuristic and strategy from the task's own
features rather than from a flag --- a goal of \code{Kw} conjuncts and a small
designated set --- and solves it at depth four:

\begin{lstlisting}
[main] Heuristic: knowledge-spread (auto, rule 'kw-only-goal')
[main] Strategy:  AO* (auto, rule 'sensing-small-designated')
[aostar] Solution found at depth 4  Expanded=32 Generated=35 Memo=6/6
[validator] OK -- 4 leaves, 6 branches checked
\end{lstlisting}

The policy is a tree, not a sequence:

\begin{lstlisting}
goto-door3_r1
  ev0: inspect3_r1
    ev0: goto-door6_r1 -> inspect6_r1     # room 3 seen blocked
    ev1: goto-door6_r1 -> inspect6_r1     # room 3 seen clear
\end{lstlisting}

Both branches continue to the second room, and they are distinct nodes rather
than a shared subtree: the executor runs one of them, and which one is decided
by an observation that has not been made when the plan is built.

\section{Where the observation comes from}

The point of the demonstration is that nothing in the stack invents the answer.
The perception node watches an occupancy grid, classifies a named region of it,
and reports what it found as the outcome of the sensing action the policy
branched on:

\begin{lstlisting}
epistemic_perception:
  ros__parameters:
    regions: ["room3", "room6"]
    room3:
      boxes: [4.25, 0.25, 11.8, 7.8]
      atom: "blocked-r3"
      atom_true_when_clear: false
      sensing_action: "inspect3_r1"
      outcome_when_clear: "e-inspect3-clear"
      outcome_when_blocked: "e-inspect3-blocked"
\end{lstlisting}

A region counts as \emph{clear} only when every cell in it has been observed and
settled, and as \emph{blocked} when any one cell is occupied; anything else is
undecided. Occupied beats unobserved, because no further looking removes an
obstacle, and unobserved beats free, because a region is clear only when all of
it is.

The grid is built by \code{slam\_toolbox} from the robot's own laser as it
drives. This is what makes the demonstration honest rather than arranged: a room
the robot has not entered is unobserved in the grid and undetermined in the
Kripke model at the same time, and those are two views of one fact. Nothing
publishes a prepared map.

\section{Execution}

The six-node epistemic bringup comes up on this domain --- domain expert,
problem expert, planner, executor, epistemic state, epistemic perception ---
and the planner answers with the policy above. The action mapping joins the two
vocabularies, \code{inspect3\_r1} on the epistemic side being
\code{(look\_into r1 door3)} in PDDL, and the executor dispatches it:

\begin{lstlisting}
[planner]    [epistemic] policy with 6 nodes, branching
[executor]   Accepted new goal
[perception] 'room3' is blocked: inspect3_r1 observed e-inspect3-blocked
[epistemic]  applied inspect3_r1 -> e-inspect3-blocked
[bt]         (look_into r1 door3) observed e-inspect3-blocked
[epistemic]  applied goto-door6_r1
[perception] 'room6' is clear: inspect6_r1 observed e-inspect6-clear
[bt]         (look_into r1 door6) observed e-inspect6-clear
[bt] [goal]  (and (Kw r1 blocked-r3) (Kw r1 blocked-r6))
Successfully finished
\end{lstlisting}

The two rooms answer differently, so the branch is exercised rather than
implied: the policy takes its blocked continuation out of room~3 and its clear
one out of room~6, and the goal holds because the agent has come to know both
propositions --- from a map it built itself while driving.

\subsection{What the run cost}

Measured against Gazebo's own poses for the robot rather than against its
odometry, which keeps integrating while a base is held against a wall and so
cannot answer the question at all:

\begin{center}
\begin{tabular}{@{}lr@{}}
\toprule
Distance driven & 27.88\,m \\
Closest approach to any obstacle & 0.307\,m \\
Robot body radius & 0.105\,m \\
Contacts & 0 \\
Recoveries from a stall & 0 \\
\bottomrule
\end{tabular}
\end{center}

\section{Five defects the demonstration found}

Building it was the test. Each of these stopped the run, and each is fixed.

\paragraph{A dangling reference in the executor} \code{EpistemicBTBuilder}
terminated with \code{SIGSEGV} whenever a policy was rendered.
\code{ExecutorNode} iterated
\code{get\_parameter("bt\_node\_plugins").as\_string\_array()} directly in a
range-\code{for}. That accessor returns a reference \emph{into} the
\code{rclcpp::Parameter} the call returned by value; the Parameter is a
temporary and dies at the end of the full expression, so the loop walked a
destroyed vector and handed \code{dlopen} a freed pointer --- a fault inside
\code{strchr}, as the loader scanned the name for a dynamic string token. It
bites only when \code{bt\_node\_plugins} is non-empty, which is exactly the
epistemic configuration and never the classical one, which is why no existing
test caught it.

\paragraph{A parameters file that could not be loaded} The shipped
\code{plansys2\_epistemic\_params.yaml} wrote \code{regions: []}. An empty list
in a parameters file has no element type, so the override reaches the node with
no value and construction throws --- the very trap the same file documents
twenty lines earlier for \code{epddl\_libraries}. Launching with
\code{epistemic\_perception:=True} therefore aborted every time.

\paragraph{Rolling-only APIs on a Humble target} Five packages used
\code{rclcpp::experimental::executors::EventsExecutor} and
\code{rclcpp::ServicesQoS()}, neither of which exists in Humble, so the fork
could not be built on the distribution its own CI names.
\code{plansys2\_core/Compat.hpp} now selects between them by asking whether the
header exists rather than what the distribution is called.

\paragraph{A refusal that was only a race} Perception reports the moment a
region resolves; the epistemic state learns the robot reached the doorway only
when the executor applies the preceding action. Observed gap: 78\,ms. The state
refused the early observation as a divergence between model and world, and
perception --- by design --- never retried a refusal. A reading the model was
about to be ready for was thrown away permanently. Perception now distinguishes
\emph{not applicable yet} from a real disagreement and retries the first for a
bounded number of grids.

\paragraph{An update applied twice} With the observation landing first, the
behaviour tree's own \code{ApplyEpistemicUpdate} then applied the same sensing
action again and was refused, because its \code{Kw} precondition no longer held
--- it had succeeded. The state now answers a re-applied action with the
outcome it already produced, and repeats no update. Every other inapplicable
action is still refused exactly as before.

\section{What a sensing action turned out to be}

The demonstration also settled a modelling question the domain alone does not
answer: what a robot should physically \emph{do} when the policy says
\code{inspect3}.

A timed pause is the obvious reading and the wrong one --- it cannot tell
success from failure, and it ends on a stopwatch rather than on knowledge. A
rotating sweep works but only as a workaround: \code{slam\_toolbox} integrates
nothing from a base that never moves, so the robot was spinning to keep its own
mapper alive, which reads as confusion rather than intent.

What the action does now is drive into the room, turn to face the region, hold
still, and end exactly when \code{(Kw r1 blocked-r3)} holds in the epistemic
state --- or fail, and say so, if the region never settles. A sensing action
exists to acquire knowledge, so knowledge is what ends it.

One consequence is worth recording, because it is a fact about sensors and not
about the model. A region is \emph{blocked} as soon as one cell in it is
occupied, and \emph{clear} only when every cell has been observed and settled.
The first is decided in an instant; the second cannot be decided at all for a
region lying in open floor, because a 3.5\,m laser returns nothing in a
7.5\,m room and free space is only traced along beams that come back. The
region a robot can settle is the region its beams cross on the way to a wall.

\section{Reproducing}

\begin{lstlisting}
# ground and solve, without ROS
plank export -d building-rooms.epddl -p instances/problem_1.epddl \
      -l intermediate.epddl -o out
epistemic_planner --task out/problem_1.json --plan out/plan_1.json --explain

# the full system
ros2 launch eplansys_rooms_demo rooms_demo_launch.py
\end{lstlisting}

The domain and instance are in \code{epddl-workspace/building-rooms}; the ROS
side, including the two action nodes, the parameters and the launch file, is in
\code{ros2\_ws/src/eplansys\_rooms\_demo}.

\end{document}
